\section{What warrants a class, a category, or an explanation?}\label{sec:assessment}

The expressive--depictive comparison establishes a descriptive distinction.
Whether either grouping supports further induction is a different question.
The same distinction applies to grammatical and interactional proposals:
an inventory can organize the observed uses accurately while leaving
its reach beyond those observations unknown. A common assessment must
recognize both the value of that description and the additional commitment
made when an empirical category is proposed.

\subsection{Bearers, grouping, contribution, and alternatives}

Four questions organize the assessment. What bears membership? What groups
the members? What does the grouping help an investigator distinguish,
infer, or explain? And what alternative would serve the same inquiry?
Each proposal also needs a scope: a language or community, relevant
conditions, and, where the claim concerns change, a period. These
questions apply to established proposals and new candidates alike.

The \term{bearer} is the object grouped. Here it can be a conventional
sign-type; a phonologist may instead group realizations, and a historian
may group changes linking earlier and later signs. A frequency of
realizations doesn't immediately establish a number of types or a rate
of change. The grammatical category of a head likewise differs from
the external distribution of its containing phrase.

The \term{grouping basis} is the evidence or relation used to assign
members: distribution, expressive contribution, ancestry, or participation
in a recognized repertoire, for example. A grouping can be sharp or
graded and can use one diagnostic or several. A unique mechanism isn't
a prerequisite for stating the proposal. Whether the proposed members
support any particular induction or explanation remains to be assessed.

The \term{contribution} identifies what the grouping accomplishes. A
constructed class can fix a comparison set. A descriptive class can
organize a recurrent relation across independently identified cases.
A category claim requires evidence that co-occurring properties support
induction beyond the application conditions. A causal claim requires
evidence about the process producing an outcome. These achievements
can accumulate, but evidence for an earlier one doesn't establish a
later one automatically.

Descriptive success can be assessed without requiring a second-property
prediction. The expressive grouping preserves a semantic relation across
standalone and integrated constructions; the expressive--depictive contrast
records different contributions concealed by the broad cover term.
Those are warranted uses of descriptive classes. Calling them descriptions
doesn't diminish their scientific value. It specifies the achievement
so that categoryhood, in the stronger sense used here, remains an
answerable further question.

An \term{alternative} can be a broader class, narrower subtypes, a
continuous dimension, or an item-by-item account. In the worked semantic
comparison, separate classes and a superordinate class with explicit
subtypes both preserve the decisive distinction. For an inquiry into
recipient responses, action-specific groupings may be more informative
than their union. The number and rank of categories are outcomes of
comparison, rather than requirements supplied by a field or inquiry.

\begin{table}[htbp]
\centering
\small
\caption{Assessing the expressive grouping without conflating its achievements}
\label{tab:inquiry-record}
\begin{tabular}{@{}>{\raggedright\arraybackslash}p{.17\textwidth}
                  >{\raggedright\arraybackslash}p{.36\textwidth}
                  >{\raggedright\arraybackslash}p{.39\textwidth}@{}}
\toprule
Question & Proposed answer & Evidential consequence \\
\midrule
Bearer & Conventional uses, including standalone and prenominal \mention{damn} &
Identify each use without requiring the target property under investigation. \\
\addlinespace
Grouping & Deployment conventionally conveys evaluative or affective stance &
Distinguish this from describing an emotion, depicting a scene, or managing interaction. \\
\addlinespace
Contribution & Preserve expressive commonality across grammatical constructions &
This is a descriptive achievement; prediction about a further diagnostic needs additional evidence. \\
\addlinespace
Alternative & An undifferentiated interjection class, explicit subtypes, or item-specific accounts &
Compare which semantic relations each preserves and what further generalization each warrants. \\
\bottomrule
\end{tabular}
\end{table}

\subsection{From proposal to category claim}

A candidate earns consideration by making these questions answerable.
Support depends on the contribution claimed. A coherent classification
can warrant continued descriptive use while remaining a prospective
category proposal. The conventional category labels inherited from the
literature aren't independent evidence that every extension considered
here has a projectible profile.

The semantic comparison gives the stages concrete content. Assigning
expressive uses by their conveyed stance fixes a class. Showing the
same contrast across grammatical constructions establishes descriptive
reach. Establishing that independently identified members share a
further interpretive property would support category standing at the
scope tested. The last claim can't be secured merely by defining the
class to contain signs already known to share every target property.

Supported category standing is consequently limited to a profile and
scope. A projection established for a subtype doesn't automatically
extend to its superordinate class. Evidence for English expressive
modifiers doesn't establish the same profile for all repair initiators
or all ideophones. These limits permit comparison and revision without
requiring one universal admission standard for every useful analytical
tool.

\subsection{Projection and identity}

A projective claim states what's been observed, what's expected, and
over which bearers and conditions. Its \term{warrant} is the evidence
supporting that expectation \citep{Goodman1955}. A precise specification
can have little warrant; a reliable association can remain causally
unexplained. The content, evidential status, and proposed ground of a
claim therefore need separate assessment
\citep[11--12]{reynolds2026projectibilityHistory}.

Suppose source contexts identify an expressive use and the target is
its interpretation under repetition. Repetition can't also supply the
case-level evidence used to admit that use. Similarly, if resistance
to determiner combination is the target, an unexpected determiner
can't by itself justify excluding the case as conversion. A change
to the identification rule creates a revised claim requiring assessment.

Consider a hypothetical adverse case for standalone stance \mention{damn}.
The initial identification uses its conventional reaction meaning and
recurrence across independently sampled situations; the grammatical
prediction concerns resistance to integration. A new occurrence retains
that reaction meaning but appears inside a larger construction. The
first response is to check the analysis of that occurrence and revise
the proposed distribution if it stands. Its inconvenience isn't evidence
for a second sign.

A split becomes justified if independent observations establish a
recurrent conventional use with a different contribution or distribution,
as the familiar prenominal pattern does. A quoted token or a productive
noun denoting an utterance is another possible analysis, but it needs
evidence for quotation or nominal reference. In the hypothetical case,
repetition of the unexpected distribution with the same reaction meaning
would count against the original restriction, even if no new sign
were recognized.

An informative restriction also needs relevant opportunities. Failure
to observe inflection says little if the sampled environments wouldn't
elicit it under a rival analysis. A grammatical test needs source and
target constructions; an interactional test needs actions and sequential
environments. The common bearer allows these claims to concern the
same use. It doesn't make their evidence interchangeable.

\subsection{Representation and causal explanation}

A category label computed deterministically from fully observed features
adds no information to those features. A representation can nevertheless
help an investigator working with incomplete or noisy observations:
it can support pooling, reduce separate assumptions, or express relevant
dependencies. Such benefits require a comparison for the specified task.
They don't establish that the represented grouping has a unique causal
ground.

Intersections receive the same treatment. A joint grouping can be
economical even though its membership follows from the component
assignments. An advantage over one restricted model wouldn't give it
priority for every inquiry. Conversely, derivability alone doesn't make
it useless. The comparison must identify what the alternatives retain
and what the investigator needs to infer.

A historical explanation can remain informative after present structure
is fully observed. If recruitment produced that structure, conditioning
on it can remove the association through which the history is recognized.
A test of additional prediction from present evidence therefore differs
from a test of the proposed history. For processing, too, conditioning
on measured frequency doesn't by itself identify repetition's total or
direct causal contribution to retrieval.

\subsection{Hill's considerations for causal claims}\label{sec:hill}

\textcite[295--299]{hill1965} offered nine viewpoints for assessing a
causal interpretation of epidemiological associations: strength,
consistency, specificity, temporality, biological gradient, plausibility,
coherence, experiment, and analogy. He rejected a mechanical checklist:
none guarantees proof, and the list isn't a set of individually mandatory
requirements. Its purpose is to examine whether a causal explanation
fits the evidence better than alternatives.

Applied to linguistic inquiry, the viewpoints concern a specified
cause and outcome, not the existence of a class. A claim that changed
prosodic packaging contributes to changed distribution needs separate
measures of both. Strength and consistency concern the observed relation;
specificity directs attention to alternative outcomes and causes;
temporality constrains the proposed direction. Plausibility and coherence
connect the claim to independently supported knowledge, while experiment
and analogy can supply further comparisons.

A linguistic counterpart of biological gradient would need a justified
intensity variable and outcome. Frequency is one possible exposure
measure, but its causal role can't simply be assumed. Counting fulfilled
viewpoints would add no evidence. And Hill's caution about mandatory
requirements doesn't permit an effect to cause its own earlier proposed
cause. Section~\ref{sec:history} applies these distinctions to the
limited history of \mention{gee}.

A recurring profile can be stable while its causal ordering remains
unknown. Historical dependence can be supported without establishing
inventory-wide recurrence. Maintenance requires evidence that a process
preserves a profile; corrective control requires a response to disturbance.
I reserve \term{homeostatic} for that stronger control claim, following
\textcite{reynolds2026notEveryStableCluster}, a narrower use than
\textcite[143]{boyd1999}. No such control claim is established here.

